% MG05a. Tras el párrafo de extensión y procedencia que sigue a
% act21:tpk:logaritmo; antes de «Orden de lectura y movimiento».
\subsubsection{Jauges de l'observable multiplicatif et normalisation du générateur}
\label{mg05:calibres}

Le choix de l'exposant discret admet une comparaison finie dans laquelle restent fixes le support APP, les conditions initiales des deux curseurs, le sélecteur tritique, le calendrier d'inversion et la règle émettrice. La variation se concentre sur l'observable du feuillet multiplicatif. Soient \(\ell_{2}\) et \(\ell_{5}\) les logarithmes du groupe \((\mathbb Z/9\mathbb Z)^\times\) de générateurs \(2\) et \(5\), respectivement, prolongés par zéro sur \(\{3,6,9\}\). Ils sont comparés à \(v_3\), évaluée sur les représentants positifs \(1,\ldots,9\), et à la lecture identité sur ces mêmes représentants.

\begin{proposition}[Normalisation du logarithme multiplicatif]
\label{mg05:normalizacion}
Les deux isomorphismes du groupe des unités modulo neuf avec
\(\mathbb Z/6\mathbb Z\), normalisés en envoyant un générateur sur
\(1\), sont \(\ell_2\) et \(\ell_5\). Ils satisfont
\[
 \ell_5(x)\equiv-\ell_2(x)\pmod6
 \qquad\bigl(x\in(\mathbb Z/9\mathbb Z)^\times\bigr).
\]
Le choix du plus petit générateur positif fixe \(\ell_2=\ell_9\),
avec les valeurs de \eqref{act21:tpk:logaritmo}.
\end{proposition}
\begin{proof}
Les puissances de \(2\) parcourent \(1,2,4,8,7,5\) et se referment à la
sixième étape. Les générateurs d'un groupe cyclique d'ordre six
correspondent aux exposants premiers avec six : \(1\) et \(5\).
Leurs représentants sont \(2\) et \(2^5\equiv5\pmod9\).
Comme \(5\equiv2^{-1}\pmod9\), exprimer une unité comme puissance de
\(5\) change le signe de son exposant en base \(2\). La
normalisation par le plus petit représentant positif sélectionne \(2\).
\end{proof}

La condition d'homomorphisme seule autorise également des applications
d'image propre. La proposition distingue donc la représentation
fidèle du groupe des unités de ses quotients. Le prolongement par zéro
détermine en outre le traitement des trois résidus non inversibles ;
le registre de provenance conserve celui d'entre eux qui a été observé.

La comparaison des quatre observables parcourt les mêmes \(104\,976\) conditions initiales et regroupe leurs émissions de six fenêtres. En désignant par \(E_g\) l'émetteur d'observable multiplicatif \(g\), on obtient le recensement d'images suivant.
\begin{table}[htbp]
\centering
\small
\begin{tabular}{@{}lrrrl@{}}
\toprule
Observable & \(|\operatorname{im}E_g|\) & Racines nulles
 & Ley aditiva & Normalisation\\
\midrule
\(\ell_2\)&468&oui&oui&plus petit générateur\\
\(\ell_5\)&468&oui&oui&generador inverso\\
\(v_3\)&144&no&oui&image unitaire nulle\\
\(\operatorname{id}\)&684&non&non&résidu positif\\
\bottomrule
\end{tabular}
\caption{Comparaison de quatre jauges sur un même transport et un même domaine de conditions initiales.}
\label{mg05:tabla-calibres}
\end{table}

Les deux logarithmes satisfont la loi additive des exposants. La valuation
\(v_3\) vaut zéro sur toutes les unités, mais prend les valeurs
\(1,1,2\) sur \(3,6,9\) ; d'où la différence entre leurs deux colonnes
de contrôle. L'identité conserve le résidu positif et ne satisfait pas la loi
d'homomorphisme vers \(\mathbb Z/6\mathbb Z\) ; par exemple, pour
l'unité \(1\), l'égalité exigée serait \(1\equiv2\pmod6\).
Le recensement de \(468\) émissions est commun aux deux logarithmes et
conserve donc l'ambiguïté d'inversion. Le critère du plus petit
générateur résout cette ambiguïté avant toute lecture régionale
de constantes. Les cardinaux du tableau décrivent l'évaluation
exhaustive de ces quatre observables sur le domaine fixé ; leur
fonction est de contrôler la normalisation opérationnelle de l'émetteur.
